\begin{lemma}
\label{lemma-geometrically-irreducible}
Let $k$ be a field.
Let $R$ be a $k$-algebra.
The following are equivalent
\begin{enumerate}
\item for every field extension $k'/k$ the
spectrum of $R \otimes_k k'$ is irreducible,
\item for every finite separable field extension $k'/k$ the
spectrum of $R \otimes_k k'$ is irreducible,
\item the spectrum of $R \otimes_k \overline{k}$ is irreducible
where $\overline{k}$ is the separable algebraic closure of $k$, and
\item the spectrum of $R \otimes_k \overline{k}$ is irreducible
where $\overline{k}$ is the algebraic closure of $k$.
\end{enumerate}
\end{lemma}

\begin{proof}
Condition (1) implies (2) directly.
Under (2), taking the identity field
extension shows that the original
$R$ has nonempty spectrum, so $R\neq0$.
Let $k_s$ be a separable algebraic
closure of $k$. Its finite
subextensions $E/k$ are separable,
are directed under compositum,
and have union $k_s$.
Every original finite list of
coefficients in $k_s$ belongs to
one such $E$, so
$R\otimes_k k_s=\bigcup_E R\otimes_k E$
through the actual injective
coefficient maps.

This tensor is nonzero by a
vector-space basis argument.
Suppose $\mathfrak q_1,\mathfrak q_2$
are minimal primes in it.
For each original finite $E$,
the map $R\otimes_k E\to R\otimes_k k_s$
is flat: its module comparison
is scalar extension by the field
$k_s/E$. Going down makes the
contractions of both primes
minimal. Indeed a smaller prime
under a contraction would lift
to a smaller prime under the
given minimal one. By (2), the
ring $R\otimes_k E$ has only
one minimal prime, so these
contractions agree. Each element
of $R\otimes_k k_s$ lies in some
one of the displayed finite
tensors; membership in either
$\mathfrak q_i$ is therefore
determined by equal contractions.
Thus $\mathfrak q_1=\mathfrak q_2$.
A nonzero ring has a minimal
prime, so this proves (3),
including the required existence.

For (3) $\Rightarrow$ (4), choose
an algebraic closure $\overline k$
containing the specified $k_s$.
The extension $\overline k/k_s$
is purely inseparable, or is the
identity in characteristic zero,
by the separable-first theorem.
The map
$R\otimes_k k_s\to(R\otimes_k k_s)\otimes_{k_s}\overline k$
is a homeomorphism on spectra
by Lemma \ref{lemma-p-ring-map}
in positive characteristic, and
is the identity in characteristic
zero. The latter tensor is
identified with $R\otimes_k\overline k$
by
$(r\otimes a)\otimes b\mapsto r\otimes ab$,
whose inverse is
$r\otimes b\mapsto(r\otimes1)\otimes b$.
This proves (4) for these original
field embeddings.

For (4) $\Rightarrow$ (1), keep
an arbitrary original extension
$L/k$. A common receiving field
must contain compatible copies
over $k$; construct it as follows.
The tensor $L\otimes_k\overline k$
is nonzero by extension of a
vector-space basis. Choose a
prime $\mathfrak a$ in it and put
$$
F=\operatorname{Frac}((L\otimes_k\overline k)/\mathfrak a).
$$
The maps $a\mapsto[a\otimes1]$
and $b\mapsto[1\otimes b]$ embed
$L$ and $\overline k$ into this
nonzero field and agree on the
original $k$ by tensor balancing.
Thus they provide the needed
commutative field diagram, not
merely unrelated abstract copies.
The scalar comparison
$$
(R\otimes_k\overline k)\otimes_{\overline k}F
\longrightarrow R\otimes_k F,
\qquad(r\otimes a)\otimes b\longmapsto r\otimes ab
$$
has inverse $r\otimes b\mapsto(r\otimes1)\otimes b$.
Its source has irreducible
spectrum by Lemma
\ref{lemma-separably-closed-irreducible},
applied over the algebraically
closed $\overline k$ to the
original $R\otimes_k\overline k$
and the field $F$.

The map $B=R\otimes_k L\to R\otimes_k F$
is scalar extension by $F/L$.
It is faithfully flat: for any
nonzero $B$-module $M$, the
comparison
$M\otimes_B(B\otimes_L F)\simeq M\otimes_L F$
sends $m\otimes(b\otimes a)$ to
$bm\otimes a$, with inverse
$m\otimes a\mapsto m\otimes(1\otimes a)$;
extension of a vector-space
basis preserves nonzeroness
and exact sequences. More
explicitly, over each prime
$\mathfrak b$ of $B$ the
nonzero ring $\kappa(\mathfrak b)\otimes_L F$
has a prime. Its inverse image
is a prime of $B\otimes_L F$
contracting to $\mathfrak b$.
Hence the spectrum map is
surjective. The continuous
surjective image of the
nonempty irreducible spectrum
just proved is irreducible,
giving (1) for the original
$R\otimes_k L$.

Two further forms follow with
the same maps. First, in (2)
it suffices to test finite
Galois extensions. Every finite
separable $E/k$ embeds in a
finite separable normal splitting
field $N/k$: take a primitive
element for $E/k$ and adjoin
all roots of its separable
minimal polynomial in an
algebraic closure. These are
finitely many algebraic roots,
all separable, so the resulting
extension is finite and
separable; each embedding
permutes the roots, so it is
normal. The field extension
$N/E$ makes
$R\otimes_k E\to R\otimes_k N$
faithfully flat by the same
module comparison above.
Irreducibility of the latter
spectrum therefore descends
to the former, proving (2).

Second, for every fixed field
extension $L/k$, the original
$R$ is geometrically irreducible
over $k$ if and only if
$R\otimes_k L$ is geometrically
irreducible over $L$.
The forward direction uses
$(R\otimes_k L)\otimes_L F\simeq R\otimes_k F$
for every field $F/L$.
For the reverse direction and
any original field $E/k$, use
the same prime-of-a-tensor
construction to place $L$ and
$E$ compatibly in a field $F$.
The hypothesis over $L$ makes
$R\otimes_k F$ irreducible,
and the faithfully flat map
$R\otimes_k E\to R\otimes_k F$
descends irreducibility as
above. Every scalar comparison
is the displayed tensor map
and its displayed inverse.
\end{proof}